\documentclass[10pt,a4paper]{article}
\usepackage{fontspec,amsmath,amssymb,graphicx}
\usepackage[margin=18mm]{geometry}
\setmainfont{Times New Roman}
\setlength{\parindent}{0pt}
\setlength{\parskip}{5pt}
\allowdisplaybreaks
\sloppy
\emergencystretch=2em
\begin{document}
\section*{VD\_1\_2\_2\_1\_04 — Lời giải bổ sung}
\subsection*{Kiến thức và trình tự giải}
\textbf{Mục tiêu.} Xác định cấp và dạng của phương trình trước khi chọn cách giải. Ở mục này, từ “bậc” trong đề được hiểu theo nghĩa cấp đạo hàm của phương trình.

\textbf{Bước 1 — Liệt kê đạo hàm và đếm cấp.} Với $\partial^{i+j}u/\partial x^i\partial y^j$, cấp là $i+j$. Vì vậy $u_x$ cấp 1, $u_{xy}$ cấp 2, $u_{txx}$ cấp 3. Hạng $u$, $e^u$, $\log u$ không chứa đạo hàm, được xem là cấp 0. $(u_y)^2$ là bình phương của một đạo hàm cấp 1, không phải đạo hàm cấp 2. Lấy giá trị lớn nhất trong các cấp vừa liệt kê.

\textbf{Bước 2 — Viết rõ phần chứa đạo hàm cấp cao nhất.} Gọi cấp đó là $m$. Tạm xem từng đạo hàm cấp $m$ như một ẩn đại số độc lập. Nếu chúng bị bình phương, nhân với nhau hoặc nằm trong hàm mũ/lượng giác, biểu thức đã cho hoàn toàn phi tuyến. Nếu chúng chỉ xuất hiện tuyến tính, chuyển sang bước 3. Khi khai triển $(f(u))_x$, phải dùng $(f(u))_x=f'(u)u_x$ trước khi phân loại.

\textbf{Bước 3 — Xét hệ số và các hạng cấp thấp.} Nếu mọi hạng chứa $u$ và đạo hàm đều tuyến tính, hệ số chỉ phụ thuộc biến độc lập, phương trình tuyến tính. Nếu chỉ phần cấp $m$ tuyến tính, hệ số của phần đó phụ thuộc biến độc lập, còn phần cấp thấp phi tuyến, phương trình nửa tuyến tính. Nếu hệ số phần cấp $m$ phụ thuộc $u$ hoặc đạo hàm cấp thấp, phương trình tựa tuyến tính. Khi ghi kết luận, chỉ rõ hạng nào quyết định phân loại.

\textbf{Bước 4 — Xét thuần nhất khi phương trình tuyến tính.} Gom toàn bộ hạng chứa ẩn về một vế, viết $Lu=f(x,y)$. Nếu $f\equiv0$, phương trình thuần nhất; nếu không, không thuần nhất. Ví dụ $u_{xy}=u$ phải viết $u_{xy}-u=0$, nên là thuần nhất.

\textbf{Bước 5 — Kiểm tra miền và yêu cầu đề.} $\log u$ đòi hỏi $u>0$ đối với nghiệm thực. Một phép biến đổi như lấy log có thể cho phương trình tương đương trên một miền, nhưng cần phân biệt dạng biểu thức in trong đề và dạng đã biến đổi. Nếu đề chỉ hỏi cấp, dừng sau bước 1 và giải thích đạo hàm quyết định cấp.
\subsection*{Lời giải áp dụng cho bài}
\textbf{Bước 1 — Xác định cấp.} $u_{xx}$ là đạo hàm cao nhất, nên $m=2$.

\textbf{Bước 2 — Xét dạng của phần cấp 2.} Đạo hàm $u_{xx}$ xuất hiện tuyến tính, nhưng hệ số là $u^2$, phụ thuộc chính ẩn cần tìm. Vì thế không thể xem phương trình là tuyến tính hoặc nửa tuyến tính theo định nghĩa trong sách.

\textbf{Kết luận.} Phương trình \textbf{tựa tuyến tính, cấp 2}. Tại $u=0$, hệ số phần chính có thể triệt tiêu; cần lưu ý sự suy biến khi giải, nhưng phân loại biểu thức chung vẫn như trên.\end{document}